% ============================================================
%  Глава 0: Введение
% ============================================================

\section*{Введение}
\addcontentsline{toc}{section}{Введение}
\label{sec:intro}

\subsection*{Мотивация}

Современный рынок биржевых и внебиржевых опционов ежедневно генерирует
колоссальный поток данных о ценах. Маркет-мейкер, управляющий книгой
опционов на S\&P~500, вынужден калибровать стохастическую модель волатильности
по реальной поверхности подразумеваемой волатильности (implied volatility, IV)
в режиме реального времени~--- за секунды после каждого изменения котировок.
Аналогичная задача возникает при управлении риском (вычисление греков),
XVA-вычислениях и стресс-тестировании портфелей.

Классические модели~--- Black--Scholes (1973)~\cite{BlackScholes1973},
Heston (1993)~\cite{Heston1993}, SABR (2002)~\cite{Hagan2002}~--- допускают
аналитические или полуаналитические решения, однако недостаточно гибки
для описания эмпирических феноменов:
\begin{itemize}
  \item \textbf{Наклон улыбки волатильности} (volatility skew) при коротких
        сроках до погашения;
  \item \textbf{Степенное поведение ATM-волатильности} вида
        $\sigma_{\mathrm{ATM}}(T) \sim C \cdot T^{H-1/2}$ с показателем
        Хёрста $H \approx 0{,}1$, характерным для дробного броуновского движения
        (Gatheral~et~al., 2018)~\cite{Gatheral2018};
  \item \textbf{Долгую память} волатильности: автокорреляция реализованной
        волатильности убывает по степенному закону, а не экспоненциально.
\end{itemize}

Модели \textit{грубой волатильности} (Rough Volatility), в частности
\textbf{Rough Heston} (El~Euch~\&~Rosenbaum, 2019)~\cite{ElEuch2019},
воспроизводят эти свойства, управляя дисперсионным процессом дробным броуновским
движением с $H < 1/2$. Однако цена этой реалистичности~--- \textbf{вычислительная
сложность}: калибровка Rough Heston классическими методами (Монте-Карло,
численное решение дробных уравнений Риккати) занимает~от~нескольких~минут до
нескольких часов, что делает их неприменимыми для реального трейдинга.

\subsection*{Цель и задачи работы}

\textbf{Цель работы}~--- разработать систему быстрой калибровки модели Rough Heston,
которая обеспечивает \textbf{латентность менее $1$ секунды} при \textbf{точности,
сопоставимой с референсным решателем}, и исследовать её применимость для задач
хеджирования и реал-тайм риск-менеджмента.

Для достижения цели решались следующие \textbf{задачи}:
\begin{enumerate}
  \item Реализовать GPU-ускоренный \textbf{прайсер Фурье--COS} с устойчивым
        ODE-решателем (exponential midpoint) для генерации датасетов поверхностей IV.
  \item Разработать и обучить \textbf{суррогатную модель} на основе
        Fourier Neural Operator (FNO) с кондиционированием по параметрам
        через Feature-wise Linear Modulation (FiLM).
  \item Построить \textbf{быстрый калибратор} на основе метода
        Ньютона--Гаусса с автоматическим дифференцированием (forward-mode AD),
        используя суррогат как оракул функции и якобиана.
  \item Провести \textbf{теоретический анализ идентифицируемости} параметров
        через матрицу Фишера (FIM) и обосновать 3D-репараметризацию.
  \item Расширить пространство параметров, включив \textbf{обучаемый показатель
        Хёрста} $H$, и оценить влияние на качество калибровки.
  \item Продемонстрировать \textbf{практическую применимость}: реал-тайм
        потоковая калибровка, авторегрессионное хеджирование по дельте.
\end{enumerate}

\subsection*{Научная новизна}

\begin{enumerate}
  \item \textbf{FiLM-FNO для прайсинга}: применение FNO с FiLM-кондиционированием
        к задаче аппроксимации поверхности IV Rough Heston достигает
        $R^2 = 0{,}9991$, $\text{MAE} = 0{,}058\%$ при инференсе $< 1\,\text{мс}$~---
        в~$10^3$-$10^4$ раз быстрее прямого прайсинга.
  \item \textbf{Автоградиентный метод Ньютона}: первое применение \texttt{jacfwd}
        (forward-mode AD через FNO) для вычисления точного якобиана
        $\partial \IV / \partial \theta$ за $3$ прохода вместо $10$ при конечных
        разностях, с квадратичной сходимостью к минимуму.
  \item \textbf{Идентифицируемость через FIM}: доказана численно разница числа
        обусловленности $\kappa(\mathbf{I}) = 1301\times$ между 6D и 3D
        параметризацией, обосновывающая фиксацию $\kappa, \theta, H$.
  \item \textbf{Обучаемый показатель Хёрста}: расширение 3D-калибратора до 4D
        с включением $H \in [0{,}04; 0{,}15]$ как свободного параметра.
\end{enumerate}

\subsection*{Практическая значимость}

Все вычисления выполняются на потребительском GPU NVIDIA RTX~3060 (6~ГБ VRAM).
Медианная латентность калибровки $549\,\text{мс}$, $p_{95} = 668\,\text{мс}$~---
достаточно для обработки тиков SPX-опционов ($\approx 1$--$2$ секунды между
тиками) в реальном времени.

\subsection*{Структура работы}

\begin{description}[leftmargin=2cm, style=nextline]
  \item[Глава~1.] \textit{Математические основы}: модель Rough Heston,
        дробное броуновское движение, лифтинговое представление, условия
        отсутствия арбитража для поверхностей IV.
  \item[Глава~2.] \textit{Прайсинг через метод Фурье--COS}: характеристическая
        функция Rough Heston, дробное уравнение Риккати, GPU-реализация,
        сходимость по $N$ факторам Бернштейна.
  \item[Глава~3.] \textit{Суррогатная модель FNO}: архитектура операторной сети,
        FiLM-кондиционирование, регуляризация на отсутствие арбитража, обучение.
  \item[Глава~4.] \textit{Калибровка}: анализ идентифицируемости, FIM,
        3D-репараметризация, метод Ньютона--Гаусса с \texttt{jacfwd}.
  \item[Глава~5.] \textit{Эксперименты}: точность прайсера, качество FNO,
        сравнение Newton vs L-BFGS, потоковая калибровка, хеджирование,
        обучаемый показатель Хёрста.
  \item[Заключение.] Выводы и перспективы.
\end{description}
